%! Introducing Statistics: What Can We Learn from Data? — Unit 1, Lesson 1.0 — Homework (Answer Key)
% GRADUAL RELEASE: this is the lesson's SCORED practice and the LAST component in the packet.
% Every lesson gets a generated homework; when DeltaMath carries the topic instead, this page
% is simply not assigned and the cover's score cell is struck through.
% Started in the closing minutes of the period and finished at home.
% \answerspace{H}{} reserves exactly H in the blank; the key passes the answer as the second
% argument, so the two files paginate identically by construction. Keep the macro, every
% \answerspace height, and every \boxguard byte-identical in the blank and the key.
% The next-lesson preview closes the packet, so it lives here rather than in ../ap_practice/.
\documentclass[10pt]{article}
\usepackage{apstats-article}
\usepackage{apstats-key}   % pulls in -boxes; do NOT also load -boxes
\newcommand{\answerspace}[2]{\par\nopagebreak\noindent\begin{minipage}[t][#1][t]{\linewidth}%
  \color{keyred}\bfseries #2\end{minipage}\par}

\begin{document}

\pageheader{Unit 1, Lesson 1.0}{Homework --- Answer Key}
% NAMESTRIP: no \namedateperiod here — mirrors the blank.
% NO teachernote here — scoring guidance lives in the lesson plan.

\vspace{0.12in}

% Authored BYTE-IDENTICALLY in homework/ and homework_key/ — this box is what keeps the
% blank and the key the same number of pages.
\boxguard[8]
\begin{remindbox}
\small
\textbf{This is your graded homework.} It is scored and due next class. You will start it in the
last minutes of the period, so ask about anything that stalls you before you leave, then finish
it on your own. Write every explanation \emph{in the context of the problem} --- that is the
standard here, and it is what earns credit on the AP exam.
\end{remindbox}

\vspace{0.10in}

\boxguard[24]
\begin{scenariobox}[The Food Truck Lot]{navy}
\small
A city keeps a record of every food truck licensed to park in one downtown lot.

\vspace{6pt}
\renewcommand{\arraystretch}{1.15}
\begin{tabularx}{\linewidth}{@{} l Y Y Y Y @{}}
\rowcolor{sky}
\textbf{Truck} & \textbf{Cuisine} & \textbf{Menu items} & \textbf{Avg.\ wait (min)} & \textbf{Cash only?} \\
\hline
Bao Down     & Chinese & 14 & 11 & No  \\
Curry Cart   & Indian  &  9 &  7 & Yes \\
El Farolito  & Mexican & 22 & 15 & No  \\
Gyro Hero    & Greek   &  6 &  4 & Yes \\
Nonna's      & Italian & 18 & 19 & No  \\
Pho Sure     & Vietnamese & 11 & 6 & No \\
\hline
\end{tabularx}
\end{scenariobox}

\vspace{0.10in}

\boxguard[18]
\begin{notesbox}{Part A --- Reading the Record}
\small
\begin{enumerate}[leftmargin=1.9em, itemsep=8pt]
  \item The individuals are \ans{the food trucks}. There are \ans{6} of them, and
        \ans{4} variables are recorded about each.
  \item ``Greek'' is a \ans{value} of the variable \ans{Cuisine}, and
        ``19'' is a \ans{value} of the variable \ans{Avg.\ wait (min)}.
  \item Write one \textbf{statistical question} the city could answer with this record, and
        state exactly what varies in it.
        \answerspace{1.6cm}{Answers vary. For example: ``How do the average waits at the trucks
        in this lot vary?'' What varies is the average wait time, which differs from one truck
        to the next (4 minutes at Gyro Hero, 19 at Nonna's). Any question naming a characteristic
        that differs across the trucks and requires reading the record earns credit.}
\end{enumerate}
\end{notesbox}

\vspace{0.10in}

\boxguard[18]
\begin{notesbox}{Part B --- Statistical or Not?}
\small
\begin{enumerate}[leftmargin=1.9em, itemsep=8pt]
  \item \textbf{Same lot, two questions.} Exactly one is statistical. Circle it, then name the
        check the other one fails and explain why it fails.

        \vspace{4pt}
        \hspace{1.0em}\emph{(i)} How many trucks are licensed for this lot?
        \hfill
        \emph{(ii)} How do the average waits at these trucks vary?
        \ans{$\leftarrow$ circle (ii)}
        \answerspace{1.8cm}{(ii) is statistical. (i) fails \textbf{check 1}: the number of
        licensed trucks is one fixed count --- 6 --- and that answer does not differ from one
        truck to the next. Needing to read the whole record is not the test; variation across
        individuals is.}
  \item A classmate says, ``A question can only be statistical if the answer is a number.''
        Give one question about this lot that shows the classmate is wrong, and explain how
        it passes both checks.
        \answerspace{1.8cm}{Answers vary. For example: ``What kind of cuisine do the trucks in
        this lot serve?'' Check 1: cuisine differs from truck to truck --- Chinese, Indian,
        Mexican, Greek. Check 2: you must read the record to answer. Both checks pass, and the
        answer is a set of category names rather than a number.}
\end{enumerate}
\end{notesbox}

\vspace{0.10in}

\boxguard[18]
\begin{notesbox}{Part C --- Context and the AP Standard}
\small
\begin{enumerate}[leftmargin=1.9em, itemsep=8pt]
  \item A sign at the entrance to the lot reads only: \textbf{``Average wait: 10 minutes.''}
        Name \textbf{two} pieces of context a customer would need before that number means
        anything, and then explain why the sign does not tell one customer how long
        \emph{they} will wait.
        \answerspace{2.2cm}{Two of: which trucks the average covers (all six, or only some);
        what time of day or day of week it was measured; how many customers were timed; whether
        it is a mean or a typical wait. The sign does not predict one customer's wait because it
        summarizes the whole lot --- actual waits range from about 4 minutes to about 19, and no
        single customer is guaranteed the group's average.}
  \item \textbf{AP-style multiple choice.} A researcher records, for each of 500 hospital
        patients, the patient's length of stay in days. Which of the following is a
        \textbf{statistical question} about this data set?
        \begin{enumerate}[label=\textbf{(\Alph*)}, leftmargin=2.2em, itemsep=3pt]
          \item How many patients were included in the study?
          \item How many days did patient \#307 stay?
          \item How do the lengths of stay vary among these patients?
          \item How many variables were recorded?
          \item On what date did the study begin?
        \end{enumerate}
        Justify your choice in one sentence by naming what varies.
        \answerspace{1.4cm}{\textbf{(C).} Length of stay differs from one patient to the next,
        so the answer varies across individuals and the data must be collected to find it.}
\end{enumerate}
\end{notesbox}

\vspace{0.10in}

% The next-lesson preview closes the packet.
\boxguard[10]
\begin{spiralbox}
\small
\textbf{Next lesson: The Language of Variation.} You now know that a variable is a column
heading and that its values differ from individual to individual. Next time we sort variables
into two kinds --- the ones that record a \emph{category} and the ones that record an
\emph{amount} --- and see why that difference decides every graph and every summary we are
allowed to use for the rest of the course.
\end{spiralbox}

\end{document}
